%\documentclass[11pt,a4paper,twoside]{tesis}
% SI NO PENSAS IMPRIMIRLO EN FORMATO LIBRO PODES USAR
\documentclass[11pt,a4paper]{tesis}

\usepackage{graphicx}
\usepackage[utf8]{inputenc}
\usepackage[spanish, activeacute]{babel}
\usepackage[left=3cm,right=3cm,bottom=3.5cm,top=3.5cm]{geometry}


\begin{document}

%%%% CARATULA
\def\titulo{Licenciado }
\def\autor{Mariano Bianchi}
\def\tituloTesis{Seguimiento de Objetos en Secuencias de Imágenes RGB-D}
\def\runtitulo{Seguimiento de Objetos en Secuencias de Imágenes RGB-D}
\def\runtitle{Object tracking using RGB-D image sequences}
\def\director{Francisco Roberto Gómez Fernández}
%\def\codirector{Master Yoda}
\def\lugar{Buenos Aires, 2014}
\newcommand{\HRule}{\rule{\linewidth}{0.2mm}}
%
\thispagestyle{empty}

\begin{center}\leavevmode

\vspace{-2cm}

\begin{tabular}{l}
\includegraphics[width=2.6cm]{logofcen.pdf}
\end{tabular}


{\large \sc Universidad de Buenos Aires

Facultad de Ciencias Exactas y Naturales

Departamento de Computaci\'on}

\vspace{6.0cm}

%\vspace{3.0cm}
%{
%\Large \color{red}
%\begin{tabular}{|p{2cm}cp{2cm}|}
%\hline
%& Pre-Final Version: \today &\\
%\hline
%\end{tabular}
%}
%\vspace{2.5cm}

{\huge\bf \tituloTesis}

\vspace{2cm}

{\large Tesis presentada para optar al t\'{\i}tulo de\\
\titulo en Ciencias de la Computaci\'on}

\vspace{2cm}

{\Large \autor}

\end{center}

\vfill

{\large

{Director: \director}

\vspace{.2cm}

{Codirector: \codirector}

\vspace{.2cm}

\lugar
}

\newpage\thispagestyle{empty}


%%%% ABSTRACTS, AGRADECIMIENTOS Y DEDICATORIA
\frontmatter
\pagestyle{empty}
\chapter*{\runtitulo}

% Acá iría el abstract en español (aprox. 200 palabras).
\noindent Actualmente las aplicaciones de métodos de seguimiento son muchas y su utilización aplica en diversos contextos, desde procesos industriales a entretenimiento. La creciente popularización de sensores RGB-D provoca un gran interés científico para aplicar nuevas técnicas de procesamiento de imágenes y adaptar otras conocidas a la enriquecida información que proveen estos sensores. Estos proveen información de profundidad en conjunto con texturas RGB de forma sincronizada y a una frecuencia de 30 cuadros por segundo, permitiendo obtener aplicaciones en tiempo real. En este trabajo presentamos un sistema de seguimiento separado en tres etapas y distintos métodos de detección y seguimiento en RGB, en profundidad y en su combinación. Utilizando datos de \textit{ground truth} obtenidos de una base de datos de escenas y objetos analizamos el funcionamiento de nuestro sistema



\bigskip

\noindent\textbf{Palabras claves:} tracking RGB-D, sistema de seguimiento, ICP, \textit{template matching}, estimación de pose, alineación, (no menos de 5)


\cleardoublepage
%\begin{center}
%\large \bf \runtitle
%\end{center}
%\vspace{1cm}
\chapter*{\runtitle}

\noindent In a galaxy far, far away, a psychopathic emperor and his most trusted servant -- a former Jedi Knight known as Darth Vader -- are ruling a universe with fear. They have built a horrifying weapon known as the Death Star, a giant battle station capable of annihilating a world in less than a second. When the Death Star's master plans are captured by the fledgling Rebel Alliance, Vader starts a pursuit of the ship carrying them. A young dissident Senator, Leia Organa, is aboard the ship \& puts the plans into a maintenance robot named R2-D2. Although she is captured, the Death Star plans cannot be found, as R2 \& his companion, a tall robot named C-3PO, have escaped to the desert world of Tatooine below. Through a series of mishaps, the robots end up in the hands of a farm boy named Luke Skywalker, who lives with his Uncle Owen \& Aunt Beru. Owen \& Beru are viciously murdered by the Empire's stormtroopers who are trying to recover the plans, and Luke \& the robots meet with former Jedi Knight Obi-Wan Kenobi to try to return the plans to Leia Organa's home, Alderaan. After contracting a pilot named Han Solo \& his Wookiee companion Chewbacca, they escape an Imperial blockade. But when they reach Alderaan's coordinates, they find it destroyed - by the Death Star. They soon find themselves caught in a tractor beam \& pulled into the Death Star. Although they rescue Leia Organa from the Death Star after a series of narrow escapes, Kenobi becomes one with the Force after being killed by his former pupil - Darth Vader. They reach the Alliance's base on Yavin's fourth moon, but the Imperials are in hot pursuit with the Death Star, and plan to annihilate the Rebel base. The Rebels must quickly find a way to eliminate the Death Star before it destroys them as it did Alderaan (aprox. 200 palabras).

\bigskip

\noindent\textbf{Keywords:} War, Rebellion, Wookie, Jedi, The Force, Empire (no menos de 5).

\cleardoublepage
\chapter*{Agradecimientos}
\begin{itemize}
    \item A mi director Francisco Roberto Gómez Fernández, Pachi para todo el mundo, por su paciencia infinita desde el día cero. Desde que nos sentamos a tomar un café para charlar sobre qué temas de imágenes podían llegar a interesarme hasta el día antes de la defensa de tesis se mostró alegre, entusiasmado y sobre todo, con muchas ganas de guiarme y ayudarme. ¡Gracias Pachi!
    \item A mis viejos, Victor y Sonia, por bancarme todos estos años tanto emocionalmente como monetariamente. Sin su apoyo esto no hubiera sido posible.
    \item A mi hermanita mayor, Paola, por acompañarme siempre y por tantos años lindos de convivencia juntos en Capital.
    \item Al bombonazo de mi novia, Natalin, por todo el aguante incondicional. Por bancarme todos estos años de carrera, algunos de los cuales me aguantó un poco loco. Gracias por recorrer este largo camino a mi lado, por tantas charlas interminables dándome aliento para no aflojar.
    \item A mi tío Willy, por haber cedido desinteresadamente su departamento para usarlo de bunker de estudio todos estos años. Claramente no hubiese sido tan fácil el camino sin haber recibido semejante ayuda. ¡Gracias infinitas!
    \item A mis amigos de la vida: Juani, ``el negro'', ``el rubio'', Fede, Luis, Nico, Gustav. Gracias por tantas juntadas, mateadas, asados, pileteadas, voleys y un largo etcétera que hicieron que los momentos libres de todos estos años no tuvieran desperdicio.
    \item A Pablo Brusco, compañero y amigo desde el inicio de la carrera. Por tantas juntadas dentro y fuera de la facultad, por estudio o por ocio. Siempre con su característica buena onda para que este largo viaje fuese un viaje de placer en primera clase. Sin él la facultad no hubiese sido tan entretenida y claramente hubiese sido un viaje mucho más largo.
    \item A Matías López y Rosenfeld, por tantas mateadas, consejos, charlas, fernets y bicicleteadas hasta la facultad juntos. Pasó de ser uno de mis primeros docentes en la carrera a un gran amigo.
    \item A todos los compañeros de cursada que me acompañaron durante toda o gran parte de la carrera y amigos que fui encontrando en la facultad en todos estos años (algunos de los cuales ya les perdí el rastro). Seguro me voy a saltear alguno pero intentaré hacer una lista extensiva: Carlos, Kevin, Thomas, Fede Pousa, Facu Carrillo, Pablo Echevarría, Dani Nuss, Julian, Agus Montero, ``Pape'', Martín, Mariano S., Herman, Caro Hadad, Carla, Fran Giménez, Mica, Nico ``el salteño'', Pablo Laciana, Viviana, Nacho Vissani, Nico Rosner
    \item A mis amigos matemáticos: Eli, Flor, Anto, Juanma. Por todos estos años de amistad y tantas alegrías compartidas.
    \item A todo el DC, desde los profesores, jtps y ayudantes hasta los admines y administrativos, por hacer del DC una casa de estudios de primerísimo nivel académico y por sobre todas las cosas, una gran familia.
\end{itemize}
 % OPCIONAL: comentar si no se quiere

%\cleardoublepage
%\hfill \textit{A mi viejo.}  % OPCIONAL: comentar si no se quiere

\cleardoublepage
\tableofcontents

\mainmatter
\pagestyle{headings}


\chapter{Introducción}

En la actualidad, las posibles aplicaciones de métodos de seguimiento o tracking son muchas y van desde el uso en la industria hasta juegos de consola. Un ejemplo de ello es la fabricación de barcos y autos mediante el uso de robots. Estas tareas se caracterizan por la necesidad de posicionar de manera precisa una herramienta sobre una pieza de trabajo. A través del uso de métodos de tracking se puede conocer la posición y pose de la pieza que se desea utilizar con respecto a la pose de la cámara y de esta forma saber cómo ubicar la herramienta necesaria para trabajar sobre la pieza en cuestión.

Otra área en donde se utiliza tracking de objetos es para la generación de estadísticas durante un partido de fútbol, tanto de jugadores como de un equipo, aunque las posibles aplicaciones en este contexto son mucho más amplias, como por ejemplo análisis de tácticas, verificación de las decisiones del árbitro, resúmenes automáticos de un partido, etc.

Actualmente existen sensores de profundidad que en conjunto con una cámara RGB pueden ser utilizados para detectar y seguir a una o más personas en tiempo real. De esta manera, mediante un sistema que procese las imágenes RGB-D de estos sensores, las personas puedan utilizar su cuerpo y sus movimientos para interactuar naturalmente con un dispositivo.

La utilización de sensores RGB-D se ha popularizado en los últimos años, cobrando un gran interés científico el estudio de aplicaciones y métodos capaces de procesar y entender la información que los mismos proveen.

La información de profundidad que nos provee un sensor RGB-D es un dato fundamental que nos posibilita encontrar la distancia de un objeto al sensor pudiendo recuperar su información 3D (tridimensional) junto a su textura RGB en tiempo real: 30 cuadros por segundo.
El video RGB-D que se obtiene provee una gran ayuda al mejoramiento y desarrollo de nuevas técnicas de procesamiento de imágenes y video ya conocidas. En particular, es de interés en esta tesis, el seguimiento de objetos en secuencias de imágenes RGB-D.

Un sistema de seguimiento se puede dividir en tres etapas bien definidas:
\begin{enumerate}
 \item Entrenamiento
 \item Detección
 \item Seguimiento cuadro a cuadro
\end{enumerate}

La etapa de entrenamiento consiste en obtener una representación del objeto al cuál se pretende seguir. Para llevarla a cabo se puede utilizar un patrón (template) ya conocido o aprenderlo de imágenes capturadas del mismo objeto. Este template luego se utiliza en la detección para ubicar la representación del objeto dentro de una imagen cualquiera. Una vez conocido el template no se requiere de una nueva ejecución del entrenamiento.

La segunda etapa, la de detección, radica en encontrar dentro de un frame del video al objeto en cuestión utilizando el método de detección deseado, valiéndose de la información registrada en la etapa de entrenamiento. Esta etapa se ejecuta, con el propósito de encontrar en la imagen el objeto a seguir, al comienzo del sistema de seguimiento y cuando el seguimiento cuadro a cuadro falla. Dado que la etapa de detección suele ser la más costosa en términos de desempeño computacional es deseable que se ejecute la menor cantidad de veces posible.

Finalmente, la tercera etapa consiste en seguir cuadro a cuadro el objeto detectado en la etapa anterior. Es decir, teniendo la ubicación del objeto en un cuadro de video se desea identificar la posición del mismo objeto en el siguiente frame. Esta etapa es la más importante ya que es la que se ejecuta en cada frame del video. La eficiencia del método de seguimiento es lo que determinará que todo el sistema de seguimiento se consiga realizar eficientemente. Si la técnica de seguimiento tiene una efectividad baja, es decir, no logra identificar la nueva posición del objeto en el siguiente cuadro, se debe volver a la etapa de detección cuyo desempeño computacional es mayor.


\section{Objetivos ¿va?}

El objetivo principal de esta tesis es la implementación, estudio y evaluación de un sistema de seguimiento de objetos en secuencias de imágenes RGB-D, con las siguientes características:
\begin{itemize}
 \item Performance Real-time: procesamiento de imágenes mayor a 10 cuadros por segundo
 \item Seguimiento de objetos tridimensionales con forma conocida previamente y de objetos aprendidos mediante una fase de entrenamiento previa
 \item Funcionamiento en sensores de profundidad de bajo costo (Kinect, XTion, etc.)
\end{itemize}

\chapter{Desarrollo}
\chapter{Experimentación}
\chapter{Discusión}
\chapter{Conclusiones}


%%%% BIBLIOGRAFIA
\backmatter
\bibliographystyle{alpha}
\bibliography{tesis}

\end{document}



